\chapter*{前言}
规范理论起源于对电磁现象的描述，其中包括超导现象；同时，它们也已成为基本粒子物理的核心内容，作为弱相互作用、强相互作用以及电磁相互作用的量子理论之出发点。近年来，经典规范理论由于其内在结构而受到拓扑学家与几何学家的广泛关注。物理学家与数学家的相关研究取得了若干重要成果，其中包括对全部“瞬子”（instantons）的发现，以及对其在量子理论中所起作用的部分理解。

本书提出了另一种研究视角。解析方法与拓扑方法的结合，能够为经典规范理论提供有价值的认识，而在这些问题上，单纯的几何方法尚未取得成功。我们尤其运用上述方法，分析了涡旋所对应的金兹堡–朗道方程以及单极子所对应的杨–米尔斯–希格斯方程的一些一般性质；对于这些方程，目前尚未得到闭式形式的解。

第一章和第二章介绍规范理论的基本背景。随后，作者讨论了寻找并刻画涡旋方程与单极子方程解的若干具体问题。第三章至第五章中所发展的方法具有较强的一般性，其适用范围显然不限于本书所研究的方程。第六章则汇集了一些分析工具，这些工具在多种不同类型的问题中都可能发挥作用。

作者感谢 D. Groisser 在研究过程中提出的诸多建议，以及他在校对方面给予的帮助；并特别感谢 Renate D\('\)Arcangelo 在紧张条件下完成的高质量排版工作。正是她的卓越付出，使得本书得以顺利出版。

本书旨在为经典规范理论的数学分析提供一个及时而系统的入门性介绍。作者希望，本书能够同时对理论物理学家和数学家有所裨益。
\begin{flushright}
阿瑟·贾菲（Arthur Jaffe）\\
克利福德·陶布斯（Clifford Taubes）\\
马萨诸塞州剑桥市\\
1980 年 8 月
\end{flushright}
